% ================================================================ chapter 1
\chapter{Introduction}
\label{ch:introduction}

\section{Background}
Generative artificial intelligence can now synthesize photorealistic talking-head
video and natural-sounding cloned speech with minimal effort and data. Face-swapping,
facial reenactment, lip-synchronization models, text-to-speech (TTS), and voice
conversion (VC) systems are publicly available, and their outputs are increasingly
difficult for humans to distinguish from genuine recordings. Because the visual and
acoustic streams of a clip can be manipulated \emph{independently}, a single video
can fall into one of four authenticity quadrants:

\begin{table}[h]
\centering
\caption{The four authenticity quadrants of a multimodal clip.}
\label{tab:quadrants}
\begin{tabular}{llll}
\toprule
\textbf{Quadrant} & \textbf{Video} & \textbf{Audio} & \textbf{Typical origin} \\
\midrule
RVRA & Real & Real & Genuine recording \\
RVFA & Real & \textbf{Fake} & Voice clone / TTS dubbed over real footage \\
FVRA & \textbf{Fake} & Real & Face swap / reenactment over real speech \\
FVFA & \textbf{Fake} & \textbf{Fake} & Fully synthetic avatar (lip-sync + TTS) \\
\bottomrule
\end{tabular}
\end{table}

\section{Problem Statement}
The overwhelming majority of published audio-visual deepfake detectors reduce this
problem to a single binary decision: the clip is either ``real'' or ``fake.'' This
formulation is forensically inadequate for three reasons. First, it provides no
\emph{attribution}: a journalist, court, or moderation system needs to know whether
the face was swapped, the voice was cloned, or both. Second, binary detectors tend
to shortcut on whichever modality is easier to classify, degrading accuracy on the
harder stream. Third, modern manipulations are increasingly \emph{partial}: only a
few seconds of audio or a short facial segment are altered, which a clip-level
binary label cannot express.

\section{Research Questions}
This thesis addresses the following research questions:
\begin{enumerate}[label=\textbf{RQ\arabic*.}]
  \item Can a single model produce independent, well-calibrated authenticity
        decisions for the visual and acoustic streams while remaining competitive
        with specialized unimodal detectors?
  \item Does explicit modeling of audio-visual synchronization improve detection of
        manipulations that break cross-modal coherence?
  \item Does disentangling modality-specific authenticity evidence from cross-modal
        consistency evidence improve generalization to unseen generators and datasets?
  \item Can the model temporally localize which segments of each stream were
        manipulated, and does this objective aid the clip-level decision?
  \item Is the detector robust to real-world degradations and fair across
        demographic subgroups?
  \item Can an embedding space pretrained on synthetically constructed quadrants
        built from pristine data alone generalize to real forgeries from unseen
        generators better than supervised training on generator outputs?
\end{enumerate}

\section{Objectives}
\begin{itemize}
  \item To design a unified forensic framework that detects, per modality, whether
        the video and the audio of a clip are real or AI-generated.
  \item To propose \textbf{DAVID-Net}, a disentangled cross-modal transformer with
        independent per-modality decision heads, a synchronization-consistency
        module, and per-modality temporal localization.
  \item To propose \textbf{Quadrant-Aware Contrastive Pretraining (QACP)}, a
        generator-free pretraining strategy for cross-generator generalization.
  \item To evaluate the framework rigorously across multiple benchmarks with
        cross-dataset, robustness, calibration, and fairness protocols.
  \item To deploy the detector as a public web application (FastAPI inference API
        on Hugging Face and a Next.js user interface).
\end{itemize}

\section{Contributions}
The main contributions of this thesis are:
\begin{enumerate}
  \item \textbf{QACP}: a self-supervised pretraining stage that constructs all four
        authenticity quadrants from pristine data only and structures three
        embedding spaces with a factorized supervised-contrastive objective,
        including an explicit \emph{real-but-mismatched} class that teaches the
        model that desynchronization is not forgery.
  \item A reformulation of audio-visual deepfake detection as a multi-task,
        per-modality attribution and localization problem.
  \item \textbf{DAVID-Net}, a disentangled dual-branch cross-modal transformer
        architecture implementing this reformulation.
  \item A generalization-first evaluation protocol with leave-one-generator-out and
        cross-dataset testing as first-class results.
  \item A complete, reproducible, publicly deployed detection system.
\end{enumerate}

\section{Scope and Limitations}
This work targets talking-head style video with a single primary subject and
speech audio, the dominant setting of current deepfake threats and benchmarks.
Detection outputs are probabilistic assessments, not legal proof of authenticity.
\todo{Refine scope after experiments.}

\section{Thesis Organization}
Chapter~\ref{ch:literature} reviews related work in unimodal and multimodal
deepfake detection. Chapter~\ref{ch:methodology} presents the proposed framework,
including DAVID-Net and QACP. Chapter~\ref{ch:setup} describes the datasets,
preprocessing pipeline, implementation, and evaluation protocol.
Chapter~\ref{ch:results} reports experimental results and ablation studies.
Chapter~\ref{ch:conclusion} concludes and outlines future work.
